\section{把 Attention 组装起来：Transformer 不只是一个公式}
\label{intro:22-transformer-architecture}

\begin{idea}[title={从“所有词都能互相查看”到完整网络}]
前面的 Attention 公式解决的是\textbf{某个位置怎样按相关度读取其他位置的信息}。但真实语言模型还需要输入 Token、位置、多个注意力头、非线性变换、稳定训练，以及预测下一个 Token 的输出层。
\end{idea}

\subsection{第一步：Token ID 转成向量，再加位置信息}

给一句话的 Token ID 序列 $(t_1,\ldots,t_T)$ 查嵌入矩阵，得到
\[
X_{\rm token}\in\R^{T\times d_{\rm model}}.
\]
若位置编码 $P$ 也具有相同形状，则可写：
\begin{equation}
X_0=X_{\rm token}+P.
\end{equation}
两张表相加是\textbf{对应位置的数字相加}。位置编码可能采用固定的正弦、可训练位置向量，或者 RoPE 这样的旋转位置方法；并不是所有 Transformer 都直接对嵌入矩阵加一个 $P$。

\subsection{第二步：由输入矩阵产生 Q、K、V}

设当前层输入为 $X$，每个 Token 是一行：
\[
Q=XW_Q,\qquad K=XW_K,\qquad V=XW_V.
\]
例如 $T=3$，每条 Query/Key 为 $2$ 维，Value 为 $4$ 维，那么
\[
Q,K:3\times2,\qquad V:3\times4.
\]
计算分数：
\[
S=\frac{QK^\top}{\sqrt{d_k}}\in\R^{3\times3},
\]
第 $i$ 行记录第 $i$ 个 Token 对三个 Key 的匹配程度；每行 Softmax 后和 $V$ 相乘，得到三个\textbf{带上下文信息}的新向量。

\subsection{第三步：为什么要多头注意力？}

同一段文本可能存在多种关系：指代、语法依赖、局部搭配等。只用一套 Q/K/V 投影，可能表达得不够灵活。多头注意力为第 $j$ 个头分别学习投影：
\[
H_j=\operatorname{Attention}(Q_j,K_j,V_j).
\]
然后将多个头的输出\textbf{在特征维度上拼接}，再经输出投影：
\begin{equation}
\boxed{\operatorname{MHA}(X)=
\operatorname{Concat}(H_1,\ldots,H_h)W_O.}
\end{equation}
比如两个头分别输出形状 $3\times2$，拼接后就是 $3\times4$；如果 $W_O:4\times4$，最终仍为 $3\times4$。

注意：\textbf{多头不是“把某条 Attention 算两遍然后平均”}，而是每个头具有可学习的不同投影。

\subsection{第四步：残差、归一化和前馈网络}

一个简化的 Pre-Norm Transformer Block 可以写成：
\begin{align}
H&=X+\operatorname{MHA}(\operatorname{Norm}(X)),\\
Y&=H+\operatorname{MLP}(\operatorname{Norm}(H)).
\end{align}
\textbf{残差}就是把原输入直接加回来；要求被加的两个矩阵形状相同。这样为深层网络的信息和梯度提供了直接路径。\textbf{归一化}帮助控制激活尺度、改善优化稳定性。实际模型可能使用 LayerNorm、RMSNorm 等不同方案。

前馈 MLP 通常\textbf{分别对每个 Token 的特征}做非线性变换；Attention 负责跨 Token 混合信息，两者角色互补。

\begin{center}
\begin{tikzpicture}[>=Latex,node distance=6mm,
 b/.style={draw=linegray,fill=paper,rounded corners=3pt,minimum height=11mm,minimum width=70mm,align=center,font=\small}]
\node[b] (e) {Token 嵌入 + 位置信息};
\node[b,below=of e] (att) {多头自注意力（MHA）};
\node[b,below=of att] (n1) {残差 + 归一化};
\node[b,below=of n1] (ff) {逐 Token 的 MLP};
\node[b,below=of ff] (n2) {残差 + 归一化};
\node[b,below=of n2] (o) {下一层 / 输出头};
\draw[->,thick,blue] (e)--(att);
\draw[->,thick,blue] (att)--(n1);
\draw[->,thick,blue] (n1)--(ff);
\draw[->,thick,blue] (ff)--(n2);
\draw[->,thick,blue] (n2)--(o);
\end{tikzpicture}
\end{center}

\subsection{为什么需要 Mask（掩码）？}

在\textbf{生成下一个 Token} 的解码器里，第 $i$ 个位置不应该偷看第 $i+1$ 个及更后面的目标 Token。因果掩码把这些位置的注意力分数设成非常小的值（实现中常为 $-\infty$），经 Softmax 后权重变成零。

例如 3 个 Token 的可见性：
\[
\begin{bmatrix}
\checkmark&\times&\times\\
\checkmark&\checkmark&\times\\
\checkmark&\checkmark&\checkmark
\end{bmatrix}.
\]
这能使\textbf{训练时的并行计算}保持“只能利用当前位置之前内容”的条件。

\subsection{原始 Transformer 与现代 Decoder-only LLM}

经典 Encoder--Decoder Transformer 包含：编码器自注意力、解码器被掩码的自注意力，以及\textbf{解码器对编码器输出的交叉注意力}。交叉注意力的 Query 来自解码器，Key 和 Value 来自编码器。

很多生成式 LLM 使用\textbf{Decoder-only} 架构：以因果自注意力和 MLP 堆叠为主，逐步预测下一个 Token。不能把所有 Transformer 都说成完全相同的结构。

\subsection{最后如何得到下一个 Token 的概率？}

对最后位置的隐藏状态 $\vect h_T$ 做输出线性投影，得到词表大小的 logits：
\[
\vect z=W_{\rm vocab}\vect h_T+\vect b_{\rm vocab}.
\]
再通过 Softmax：
\[
p_i=\frac{e^{z_i}}{\sum_j e^{z_j}},
\]
得到下一 Token 的概率分布。生成时按一定策略选择 Token，再把它加入上下文，重复计算。

\begin{conclusion}
Transformer = \textbf{Token 表示 + 位置信息 + 注意力 + MLP + 残差和归一化}，再按任务选择编码器、解码器、Mask 和输出头。Attention 是核心组件，但不是模型的全部。
\end{conclusion}
